\documentclass[11pt,a4paper]{article}

% Packages
\usepackage[utf8]{inputenc}
\usepackage[T1]{fontenc}
\usepackage[margin=2.5cm]{geometry}
\usepackage{graphicx}
\usepackage{hyperref}
\usepackage{enumitem}
\usepackage{titlesec}
\usepackage{parskip}
\usepackage{biblatex}

% Bibliography file
\addbibresource{bibliography.bib}

% Section formatting
\titleformat{\section}
  {\normalfont\Large\bfseries}{\thesection.}{1em}{}
\titleformat{\subsection}
  {\normalfont\large\bfseries}{\thesubsection}{1em}{}

% Hyperref setup
\hypersetup{
    colorlinks=true,
    linkcolor=black,
    filecolor=magenta,
    urlcolor=blue,
    citecolor=blue,
}

\begin{document}

% Title
\begin{center}
    \LARGE\textbf{Project Specification} \\
    \vspace{0.5cm}
    \large DD2375 Project Course in HPC \\
    \vspace{0.3cm}
    \normalsize Paul Mayer \\
    \today
\end{center}

\vspace{1cm}

\section{PROJECT INFORMATION}

\subsection*{Preliminary Title}
% Enter your project title here
Adapting Jaxley for Efficient Parameter Identification in AdEx Neuron Models: A JIT-Compiled, GPU-Accelerated Approach using Automatic Differentiation

\subsection*{Keywords}
% List 3-5 keywords relevant to your project
JIT-Compilation, Automatic Gradient Descent, GPU-Acceleration, Integrate-and-Fire Model, Simplified Neuron Model

\section{BACKGROUND \& OBJECTIVE}
A central challenge in neuroscience has been to 'identify the parameters of detailed biophysical models, such that they match physiological measurements at scale' \cite{Deistler2024} in reasonable simulation runtimes.
The zoo of optimisation techniques is large, spanning from evolutionary algorithms \cite{VanGeit2008} to gradient descent \cite{Jones2024}, however, deciding which algorithm to use may often depend on the underlying properties of the model.
\verb|JAXLEY|, a newly developed framework, uses automatic differentiation to quickly parameters to experimental data.
It shows very promising results, matching or outperforming other state-of-the-art (e.g. genetic algorithms) in performance for small and large Neurone models respectively \cite{Deistler2024}.

Neurone models, like models in all sciences, should explain observations and make predictions. However, the level of abstraction of models vary a lot depending on underlying assumptions and required parameters.
Therefore the choice of model is highly dependent on the problem at hand \cite{Levenstein2023}.
The level of abstraction not only has direct impact on the performance metrics of the simulation but also create difficulties when analysing larger and more complex networks \cite{Gerstner2014}. 
Therefore, Simplified Neurone Models like the Adaptive Exponential Integrate-And-Fire \cite{Naud2008} model are still used today \cite{Tesler2024}.

It remains unclear if adapting a Simplified Neurone Model  (e.g. AdEx or Leaky integrate-and-fire) to allow for automatic differentiation (e.g. by introducing surrogate gradients) performance gains can be achieved.
This project will try to answer this question.

\section{RESEARCH QUESTION \& METHOD}

\subsection*{Research Question}
\textbf{Question:} How could Jaxley be used to increase performance for parameter identification of adaptive exponential integrate-and-fire neural models.\\
\textbf{Hypothesis:} Making the differential equation continuous will effect the shape of the spike and therefore also the spiking behaviour.
The main challenge will be to keep the equation behaviour as close to the original equation, preferably not changing the parameter behaviour compared to the original model.
However, this will be very unlikey, I do expect deviations.
Runtime improvements should be expected, especially when making use of accelerators due to 'access' to better optimisers.

\subsection*{Objectives}
\begin{enumerate}
\item Provide discontinuous implementation of AdEx Model using Jaxley. That means: reimplemented the AdEx differential equations in the Jaxley Framework
      and verify the behaviour against other AdEx implementations.
      The discontinuity is the result of a reset, once a potential threshold is reached.
\item Provide surrogate gradient implementation in Jaxley that reproduces results of other AdEx Model implementations or make equations differentiable.
\item Compare performance of discontinuous and continuous model.
\item If time, compare performance of discontinuous and continuous model using GPU accelerators.
\item If time, compare runtime performance to other state-of-the-art parameter optimisers (e.g. using genetic algorithms) of AdEx Neurone Model
\end{enumerate}

\subsection*{Tasks}
% Describe the tasks necessary to reach the objectives
% For each task, describe the challenges it involves
\begin{enumerate}
    \item Provide discontinuous implementation of AdEx Model using Jaxley. That means: reimplemented the AdEx differential equations in the Jaxley Framework
      and verify the behaviour against other AdEx implementations.
      The discontinuity is the result of a reset, once a potential threshold is reached.

    \textbf{Challenges}:
    \begin{itemize}
        \item Documentation of Jaxley framework (very new and cutting edge research software) might be immature.
        \item Correctness needs to be verified with other simulations / experiments, can be tricky to find applicable implementations.
    \end{itemize}

    \item Allow model to be automatically differentiated, e.g. surrogate gradient (using splines?).

    \textbf{Challenges}:
    \begin{itemize}
        \item Changes to model might decrease it's precision.
        \item I might not be able to come up with a reasonable idea within the project timeline on how to make the model differentiable.
        \item I might struggle to implement this within the given time constraints.
    \end{itemize}
    \item Compare performance of differentiable and not differentiable model.

    \textbf{Challenges}:
      \begin{itemize}
        \item I think this should be pretty straight forward.
      \end{itemize}
    \item If time, compare performance using GPU accellerators.

    \textbf{Challenges}:
      \begin{itemize}
      \item Finding suitable testing environment that has GPUs available (I think Google Colab should work though)
      \end{itemize}
    \item If time, compare performance with other (state-of-the-art) AdEx implementation(s).

    \textbf{Challenges}:
    \begin{itemize}
    \item Finding other implementation might be challenging.
    \item Running said implementations in same computational environment might lead to complications.
    \item Validation against real experimental data might fail.
    \item Also, I suspect that there will just not be enough time within this course to finish this, e.g. getting other AdEx implementations to run so I can compare them.
    \end{itemize}
\end{enumerate}

\subsection*{Method}
The main contributions will be the following:

I will develop a differentiable version of the AdEx integrate-and-fire model.
This is needed because Jaxley performs automatic differentiation to do parameter optimisations.
Validate against experimental data, depending on the specific neurone type / network (e,g.: striatal projection neurons might be applicable).
To compare and quantify the runtime performance of differentiable and non-differentiable implementation, I plan to run a benchmark on given implementation. I will specify this after conducting the necessary literature review. If there is time, I want to benchmark using GPUs as well.
I aim to compare this to other state-of-the-art approaches as well.
For this I would need to measure the runtime in the same computational environment, which I think will probably a Google Colab environment.
This should indicate if exploiting the gradient information can increase performance of parameter identification.

\section{TIME PLAN}

% Provide a project timeline, specifying main tasks and time allocated
% Include milestones (time of achievement of intermediate goals)

\begin{table}[h]
\centering
\begin{tabular}{|l|l|l|}
\hline
\textbf{Task/Milestone} & \textbf{Time Period} & \textbf{Duration} \\
\hline
Literature review & Week 46 & 1 weeks \\
\hline
Development \& verification of non-contiuous equations & Week 47 & 1 week \\
\hline
Development \& verification of contiuous equations & Week 48-49 & 2 weeks \\
\hline
Milestone 1: Working implementation obtained & End of Week 49 & - \\
\hline
\hline
Performance testing & Week 50 & 1 weeks \\
\hline
Milestone 2: Results obtained & End of Week 50 & - \\
\hline
\hline
Analysis \& writing \& presentation preparation & Week 51-1 & 2 weeks \\
\hline
Final revision & Week 2 & 1 week \\
\hline
Milestone 3: Project completion & End of Week 2 & - \\
\hline
\end{tabular}
\caption{Project timeline with tasks and milestones}
\end{table}

\section{REFERENCES}
\printbibliography[heading=none]

\section{AI Disclosure}
\begin{itemize}
\item Creating a LaTeX template for this submission.
\begin{itemize}
\item Tools/models used (names and versions).\\
  Claude Code v2.0.31\\
  Model: Sonnet 4.5
\item Task: I asked claude to create a latex template for this submission that I can fill with my project idea.
\item Prompt: Could you create a latex template that I can use to write a project specification. It should follow this templat: [[template provided in 1.4 - Project Specification Template]]
\item Validation: Compiled the .tex file and removed all the prefilled data. Replaced the bibliography file with a file created by zotreo.
\end{itemize}
\end{itemize}

\end{document}
